\section{Reproduction}\label{app:repro}

Everything in this paper is produced by the Python package \texttt{collusim} (version 0.1.0; numpy and numba are its only runtime dependencies) and three scripts that sit on top of it. The repository is at \url{https://github.com/pradeep221b/algorithmic-collusion-auction}. To reproduce every table and figure, clone the repository and run, in order,
\begin{verbatim}
pip install -e ".[figures,test]"
pytest
python experiments.py
python indicators.py
python phase1.py <sweep|reinject|residual|payasbid|occupation|
                  occupation_mem1|occupation_mem1_long|fig>
\end{verbatim}
\texttt{pytest} runs the 31 tests (clearing rule, Q-update, indicators on hand-built policies, package surface). \texttt{experiments.py} runs the price-level study of Section~\ref{sec:results} and the ablations of Table~\ref{tab:m4}. \texttt{indicators.py} runs the punishment, ablation and temptation verdicts. \texttt{phase1.py} takes one subcommand per experiment of Section~\ref{sec:phase1}; \texttt{occupation} is the memory-0 occupation run, \texttt{occupation\_mem1} the memory-1 run at the same five noise levels, \texttt{occupation\_mem1\_long} the single 30M-round point at $\varepsilon = 10^{-4}$, and \texttt{fig} draws Figure~\ref{fig:phase1} from the stored results.

The training loop is compiled with numba and runs at about 0.4 s per million rounds per seed. Measured at 100 seeds, one process per job, on a 12-thread laptop with no GPU: the six base configurations take about 9 minutes in total; the ten cells of the $\beta$ sweep about 27 minutes; the two pay-as-bid runs about 4 minutes; the re-injection runs about 6 minutes; the Bellman residual about 2 minutes; the indicator battery on all six configurations about 3 minutes; each occupation curve about 10 minutes and the single 30M-round point about 20 minutes. Run sequentially the whole study takes about two hours; with the jobs overlapped, under an hour and a half. The environment used was Python 3.14 with numpy 2.x.

Every number in the paper is stored per seed, in the \texttt{results} directory. The file \texttt{<config>.json} holds, for each seed, the frozen-play $\Dl$, the convergence round and the last-round profile; the file \texttt{indicators\_\allowbreak<config>.json} holds every punishment, ablation and temptation verdict; the file \texttt{phase1\_\allowbreak<name>.json} holds the Phase-1 results. The final Q-tables are saved as \texttt{.npz} files next to the JSON summaries and are not committed to the repository; the re-injection and residual experiments read them, so those two subcommands need the base runs to have been executed first. The re-injection restarts draw from a separate seed stream, so the new exploration draws are independent of the original training draws.

One naming convention differs between the code and the paper. The code indexes firms 0, 1, 2; the paper calls them firms 1, 2, 3 throughout. Firm 1 in the paper is agent 0 in the code, and the ``deviator = firm 1'' rows of Section~\ref{sec:results} are the code's default \texttt{deviator=0}.

\section{Parameters and procedures}\label{app:params}

\paragraph{Price grid.} The 15 admissible bids are the points of \texttt{linspace(10, 100, 15)}, so rung $r$ has price $p_r = 10 + (r-1)\cdot 90/14$ \texteuro/MWh and one rung is \eur{6.43}. Table~\ref{tab:app-grid} lists them. In the body of the paper bid profiles are written with prices rounded to the nearest euro, so that (10,\,16,\,16) means rungs 1, 2, 2.

\begin{table}[htbp]
\centering
\caption{The bid grid. Notes: prices in \texteuro/MWh; the marginal cost is \eur{10} and the price cap \eur{100}.}
\label{tab:app-grid}
\begin{tabular}{rr@{\qquad}rr@{\qquad}rr}
\toprule
$r$ & $p_r$ & $r$ & $p_r$ & $r$ & $p_r$ \\
\midrule
1 & 10.00 & 6 & 42.14 & 11 & 74.29 \\
2 & 16.43 & 7 & 48.57 & 12 & 80.71 \\
3 & 22.86 & 8 & 55.00 & 13 & 87.14 \\
4 & 29.29 & 9 & 61.43 & 14 & 93.57 \\
5 & 35.71 & 10 & 67.86 & 15 & 100.00 \\
\bottomrule
\end{tabular}
\end{table}

\paragraph{Learner parameters.} Learning rate $\alpha = 0.15$, discount $\gamma = 0.95$ (or 0 in the ablation), exploration $\varepsilon_t = e^{-\beta t}$ with $\beta = 10^{-5}$, so that $\varepsilon = 0.37$ at $10^5$ rounds, $0.05$ at $3\times 10^5$ and $4.5\times 10^{-5}$ at $10^6$. With memory 1 the state is last round's full bid profile, $15^3 = 3{,}375$ states; with memory 0 there is one state. The table is initialised at $Q_0[s,a] = \mathbb{E}[\text{profit}(a,\ \text{rivals uniform})]/(1-\gamma)$. Horizons are 1M and 5M rounds and every run goes to the full horizon. These values were fixed in the specification before any result was seen and were never tuned.

\paragraph{Convergence.} A run is recorded as converged when no agent's greedy action, in any state, changed during the last 100,000 rounds. The flag is recorded and reported; it never stops a run.

\paragraph{Frozen play and $\Dl$.} Learning and exploration are switched off, the agents play greedily from the final state for 1,000 rounds, the first 100 rounds are dropped, and $\Dl = (\bar p - 10)/90$ is computed from the mean clearing price $\bar p$ of the remaining 900.

\paragraph{Punishment procedure.} Learning stays frozen. Greedy play first settles for 200 warm-up rounds. Six pre-deviation rounds are recorded. In the next round the deviator is forced to bid \eur{10}; the rivals play their greedy actions. Twenty-four post-deviation rounds of greedy play follow. A counterfactual path is generated from the same state with no forced bid. The verdict is read off the rivals' bids relative to the counterfactual in two windows, rounds 1--12 and 13--24 after the deviation, with a threshold of half a rung, $\eur{3.21}$: \emph{punish--forgive} if the rivals' mean bid is at least \eur{3.21} below the counterfactual in rounds 1--12 and back within \eur{3.21} in rounds 13--24; \emph{punish, no recovery} if below in both windows; \emph{no reaction} if the rivals do not move. A case is a \emph{no-op} when the deviator was already bidding \eur{10}, so that the forced bid changes nothing. Twelve-round windows were chosen because 12 is divisible by the cycle lengths 1 to 4 and 6; periods 5 and 7 are not covered, which affects 9 seeds in the long run.

\paragraph{Temptation.} At every state on the converged cycle and for every firm, the best one-shot deviation profit, with rivals held at their greedy actions, minus the realised profit. A seed has ``unclaimed temptation'' if any such gap is positive.

\paragraph{Bellman residual.} For every temptation on the converged cycle, indexed by firm $i$, state $s$, played action $a_i$ and best unplayed deviation $a_{\mathrm{dev}}$, let $a_{-i}(s)$ denote the rivals' greedy actions at $s$ and $s'(a)$ the state that follows when firm $i$ plays $a$ against them. Define the one-step target and residual of an action $a$ as
\begin{align*}
\mathrm{target}_i(s,a) &= \pi_i\big(a,\ a_{-i}(s)\big) + \gamma \max_{a'} Q_i\big(s'(a),\ a'\big), \\
\mathrm{resid}_i(s,a) &= \mathrm{target}_i(s,a) - Q_i(s,a),
\end{align*}
evaluated with the current, frozen table. The residual on the played action is a check: it must be near zero on a converged cycle, and it is. The residual on $a_{\mathrm{dev}}$ measures how far the stored deviation value is from what one fresh update would make it. A backup \emph{prefers the deviation} when $\mathrm{target}_i(s,a_{\mathrm{dev}}) > Q_i(s,a_i)$, that is, when a single update of the deviation entry would already flip the greedy choice. The share reported in Table~\ref{tab:residual} is the fraction of a seed's temptations at which this holds, averaged over seeds. Its limitation is that $Q_i(s'(a),\cdot)$ is itself taken from the same table and may be stale; a $k$-step rollout would be stricter and was not done.

\section{Additional results}\label{app:extra}

\paragraph{Punishment by deviator.} Table~\ref{tab:m3} reports the long-run verdicts with firm 1 as the forced deviator. Table~\ref{tab:app-deviator} repeats the split for firms 2 and 3. The shares are similar in aggregate, but the verdict for a given seed changes with the identity of the deviator in 61 of 100 seeds. In 47\% of seeds at least one of the three deviators produces a punish--forgive verdict.

\begin{table}[htbp]
\centering
\caption{Punishment verdicts in the long run (5M rounds, memory 1, 100 seeds) by identity of the forced deviator. Notes: percentages of seeds; the firm-1 row is the one in Table~\ref{tab:m3}; 2\% of firm-1 cases are no-ops (the deviator was already at \eur{10}).}
\label{tab:app-deviator}
\begin{tabular}{lrrr}
\toprule
Deviator & punish--forgive & punish, no recovery & no reaction \\
\midrule
Firm 1 & 21 & 70 & 9 \\
Firm 2 & 20 & 69 & 11 \\
Firm 3 & 17 & 73 & 10 \\
\bottomrule
\end{tabular}
\end{table}

\paragraph{Mean paths around the deviation.} Table~\ref{tab:app-path} gives the mean clearing price and the rivals' mean bid gap to the counterfactual in the two post-deviation windows, for all 100 long-run seeds and for the two main verdict subsets. Averaged over all seeds the rivals' bids sit 16.7 below the counterfactual in rounds 1--12 and 13.5 below in rounds 13--24, and the clearing price falls from \eur{66.0} before the deviation to \eur{45.1} and then \eur{49.4}, against a counterfactual of \eur{65.9} in both windows. The forgive subset falls to \eur{49.9} in the first window and returns almost exactly to its counterfactual in the second (gap 0.2, price \eur{63.2} against \eur{65.4}); the no-recovery subset stays about 20 below it and its price remains in the low forties.

\begin{table}[htbp]
\centering
\caption{Mean price paths around a forced one-round undercut, long run (5M rounds, memory 1, deviator firm 1). Notes: ``rivals' cut'' is how far the rivals' mean bid sits below its counterfactual value, in \texteuro/MWh; prices are mean clearing prices in \texteuro/MWh; the counterfactual is the same rounds with no forced bid.}
\label{tab:app-path}
\footnotesize
\setlength{\tabcolsep}{4pt}
\begin{tabular}{lrrrrrr}
\toprule
 & & \multicolumn{2}{c}{rivals' cut} & \multicolumn{3}{c}{clearing price} \\
\cmidrule(lr){3-4}\cmidrule(lr){5-7}
Seeds & $n$ & rounds 1--12 & rounds 13--24 & rounds 1--12 & rounds 13--24 & counterfactual \\
\midrule
All & 100 & 16.7 & 13.5 & 45.1 & 49.4 & 65.9 \\
Punish--forgive & 21 & 10.9 & 0.2 & 49.9 & 63.2 & 65.4 \\
Punish, no recovery & 70 & 20.5 & 19.1 & 41.8 & 43.8 & 66.2 \\
\bottomrule
\end{tabular}
\end{table}

\paragraph{Converged profiles.} Table~\ref{tab:app-profiles} lists the most common last-round bid profiles on frozen play for the $\gamma = 0$ long run and the two pay-as-bid runs. The $\gamma = 0$ learner ends on long irregular cycles: only 20 of 100 seeds have period 4 or less and periods run up to 25, and its most common profiles are one-shot Nash profiles or near them, with one firm at cost and two tied one or more rungs above. The pay-as-bid runs end at symmetric ties, each firm selling a third of demand at its own bid, at a higher level without memory than with it.

\begin{table}[htbp]
\centering
\caption{Most common last-round bid profiles on frozen play, 5M rounds, 100 seeds per run. Notes: bids rounded to the nearest euro; count is the number of seeds ending at that profile; permutations of an asymmetric profile are pooled.}
\label{tab:app-profiles}
\begin{tabular}{llr}
\toprule
Run & Profile & Seeds \\
\midrule
$\gamma = 0$, memory 1 & (10, 16, 16) & 15 \\
 & (10, 36, 36) & 10 \\
 & (10, 42, 42) & 10 \\
 & (10, 29, 29) & 9 \\
 & (16, 42, 42) & 7 \\
 & (10, 55, 55) & 7 \\
\midrule
Pay-as-bid, memory 0 & (68, 68, 68) & 16 \\
 & (94, 94, 94) & 12 \\
 & (81, 81, 81) & 9 \\
 & (87, 87, 87) & 8 \\
\midrule
Pay-as-bid, memory 1 & (55, 55, 55) & 13 \\
 & (49, 49, 55) & 9 \\
 & (61, 61, 61) & 8 \\
 & (49, 49, 49) & 8 \\
\bottomrule
\end{tabular}
\end{table}

\paragraph{Block-to-block moves in the occupation runs.} In the constant-$\varepsilon$ runs of Table~\ref{tab:occupation} the price is averaged in blocks, and a move is counted when the block mean changes by more than \eur{5}. The two learners move differently. The memory-0 learner makes large jumps: about \eur{28} in either direction at $\varepsilon \le 0.03$ and about \eur{22} at $\varepsilon = 0.1$, which is the signature of the walk down the grid and the single jump back up described in Section~\ref{sec:discussion}. The memory-1 learner makes small ones: about \eur{10} at $\varepsilon \le 0.003$, about \eur{8} at $\varepsilon = 0.01$, about \eur{4}--\eur{5} at $\varepsilon = 0.03$, and none at all at $\varepsilon = 0.1$, where its time-average $\Dl$ is 0.336 and no block exceeds $\eur{55}$.
